Let $\mathcal{S}=\{\SeqS_1,\dots\SeqS_n\}$ be a set of normalised
scenarios over $\Sigma_{\mathcal{S}}$, such that, for every $1\leq
i\leq n$, $\SeqS_i=(\E_i,\emptyset)$, where $\E_i$
%over $\Sigma_{\mathcal{S}}$, whose members
contains all the events from $\Sigma_{\mathcal{S}}$. Moreover, let
$\mathcal{S}$ be without gaps. We are interested
in obtaining an expression that realises $\mathcal{S}$.
%First, we show a simple example.
%---------------
\begin{figure}[t]
\begin{center}
\adjustbox{valign=b}
 {\begin{minipage}{0.41\linewidth}
 \begin{center}
 %               \adjustbox{valign=b}
%{\begin{minipage}{0.49\linewidth}
  \begin{tikzpicture}
    \tikzset{vertex/.style = {shape=circle,draw,minimum size=1em}}
    \tikzset{vertex3/.style = {shape=circle,draw,minimum size=1.5em}}
    \tikzset{vertex2/.style = {shape=rectangle,draw=none,minimum size=0.5em}}
    \tikzset{edge/.style = {-,> = latex'}}

    % vertices

    \node[vertex2] (b) at  (0,0) {$\mathtt{start}$};

    \node[vertex2] (x) at  (-1.5,-1.4) {$\mathtt{finish_1}$};
\node[vertex2] (y) at  (0,-1.4) {$\mathtt{finish_2}$};
\node[vertex2] (z) at  (1.5,-1.4) {$\mathtt{flag}$};
 \draw[edge] (b) to (x);
    \draw[edge] (b) to (y);
    \draw[edge] (b) to (z);
      \end{tikzpicture}
     \caption{Maximal chains (Example \ref{ex-one-same-flag})}
     \label{fig-chains-1}
      \vspace{-3mm}
     \end{center}
    \end{minipage}}
    \adjustbox{valign=b}
        {\begin{minipage}{0.58\linewidth}
         \begin{center}
 %               \adjustbox{valign=b}
%{\begin{minipage}{0.49\linewidth}
 \adjustbox{valign=b}
        {\begin{minipage}{0.49\linewidth}
  \begin{tikzpicture}
    \tikzset{vertex/.style = {shape=circle,draw,minimum size=1em}}
    \tikzset{vertex3/.style = {shape=circle,draw,minimum size=1.5em}}
    \tikzset{vertex2/.style = {shape=rectangle,draw=none,minimum size=0.5em}}
    \tikzset{edge/.style = {-,> = latex'}}

    % vertices

    \node[vertex2] (b) at  (0.2,-0.2) {$\mathtt{start}$};

    \node[vertex2] (x) at  (0.2,-1) {$\mathtt{finish_1}$};
\node[vertex2] (y) at  (-1,-1.7) {$\mathtt{finish_2}$};
\node[vertex2] (z) at  (1,-1.7) {$\mathtt{flag}$};
 \draw[edge] (b) to (x);
 %   \draw[edge] (b) to (y);
    \draw[edge] (x) to (z);
     \draw[edge] (x) to (y);
      \end{tikzpicture}
       \end{minipage}}
        \adjustbox{valign=b}
        {\begin{minipage}{0.49\linewidth}
  \begin{tikzpicture}
    \tikzset{vertex/.style = {shape=circle,draw,minimum size=1em}}
    \tikzset{vertex3/.style = {shape=circle,draw,minimum size=1.5em}}
    \tikzset{vertex2/.style = {shape=rectangle,draw=none,minimum size=0.5em}}
    \tikzset{edge/.style = {-,> = latex'}}

    % vertices

    \node[vertex2] (b) at  (0.2,-0.2) {$\mathtt{start}$};

    \node[vertex2] (x) at  (0.2,-1) {$\mathtt{finish_2}$};
\node[vertex2] (y) at  (-1,-1.7) {$\mathtt{finish_1}$};
\node[vertex2] (z) at  (1,-1.7) {$\mathtt{flag}$};
 \draw[edge] (b) to (x);
  %  \draw[edge] (b) to (y);
    \draw[edge] (x) to (z);
     \draw[edge] (x) to (y);
      \end{tikzpicture}
       \end{minipage}}
     \caption{Maximal chains for $\ParS{\mathcal{S}_1}$,
    $\ParS{\mathcal{S}_2}$(Example \ref{ex-one-same-flag})}
     \label{fig-chains-2}
     \vspace{-3mm}
     \end{center}
    \end{minipage}}
  \end{center}
  \end{figure}
%-----------------------
\vspace{-1em}
\begin{example}
\label{ex-one-same-flag}
%Consider the set $\mathcal{S}=\{\SeqS_1,\SeqS_2,\SeqS_3,\SeqS_4\}$, of
%scenarios, where
Let $\mathcal{S}=\{\SeqS_1,\SeqS_2,\SeqS_3,\SeqS_4\}$, where
%\noindent
$\SeqS_1=(\mathtt{start}~ \mathtt{finish_1}~ \mathtt{finish_2}~ \mathtt{flag},\emptyset)$,
$\SeqS_2=(\mathtt{start}~ \mathtt{finish_1}~  \mathtt{flag}~ \mathtt{finish_2},\emptyset)$,
$\SeqS_3=(\mathtt{start}~ \mathtt{finish_2}~ \mathtt{finish_1}~ \mathtt{flag},\emptyset)$
and
$\SeqS_4=(\mathtt{start}~ \mathtt{finish_2}~ \mathtt{flag}~ \mathtt{finish_1},\emptyset)$.
The set $\mathcal{S}$ cannot be realised by a single DTS, as will be explained
below.
Unlike in Example \ref{ex-one-flag},
$\mathcal{S}$ cannot be divided according to
sets of events: each member contains the same set of events.
So we need another criterion for partitioning $\mathcal{S}$.
\end{example}
%-----------------
\begin{definition}
\label{def-match}
Two (sequential) scenarios, $\SeqS$ and $p$, \emph{match} if
$\ESeq{\SeqS}=\ESeq{p}$.
%If $\SeqS$ and $p$ match then we write $match(\SeqS,p)$.
\end{definition}
%----------------
\begin{definition}
\label{def-DR-Proj}
Let $\mathcal{S}$ be a finite set of scenarios and $\DR{\mathcal{S}}$ be
its distance relation. The set of projections of $\DR{\mathcal{S}}$
(see \Sec{sec-dts})
is
denoted by $\Proj{\mathcal{S}}$.
\end{definition}
%------------------
\begin{definition}
\label{def-order-closed}
Let $\mathcal{S}$ be a finite set of scenarios.
We say $\mathcal{S}$ is \emph{order-closed} iff
every member of $\Proj{\mathcal{S}}$  matches one scenario in
$\mathcal{S}$ and vice versa.
\end{definition}
\noindent
Intuitively, if $\mathcal{S}$ is order-closed, then extending it with
a scenario that has a new sequence of events will change
$\ParS{\mathcal{S}}$.
\begin{obs}
\label{obs-order-closed-1}
Let $\mathcal{S}$ be a finite set of scenarios.
%such that all its members have only default constraints.
If $|\mathcal{S}|=|\Proj{\mathcal{S}}|$, then
$\mathcal{S}$ is order-closed.
\end{obs}
\noindent
As mentioned at the beginning of \Sec{sec-realize}, $\mathcal{S}$ can
be realised as a single DTS only if
$\mathcal{S}=\Proj{\mathcal{S}}$. This leads to the following observation:
\begin{obs}
\label{obs-order-proj-2}
Let $\mathcal{S}$ be a finite set of scenarios.
If $\mathcal{S}$
%A set of scenarios that
is not order-closed, then it cannot be realised by a single DTS.
\end{obs}
\noindent
In Example \ref{ex-one-same-flag}, $\ParS{\mathcal{S}}$
is represented by the maximal chains $\mathtt{start}-\mathtt{finish_1}$,
$\mathtt{start}-\mathtt{finish_2}$ and
$\mathtt{start}-\mathtt{flag}$ (see \Fig{fig-chains-1}).
The distance function of $\mathcal{S}$, $\DF{\mathcal{S}}$, has only
default constraints.
The set of projections of
$\DR{\mathcal{S}}=(\ParS{\mathcal{S}},\DF{\mathcal{S}})$, includes
four projections that match the
scenarios in $\mathcal{S}$.
Additionally, $\Proj{\mathcal{S}}$ includes two
projections whose sequences of events are
$\mathtt{start}~ \mathtt{flag}~ \mathtt{finish_1}~ \mathtt{finish_2}$
and
$\mathtt{start}~ \mathtt{flag}~ \mathtt{finish_2}~ \mathtt{finish_1}$.
So, by \Def{def-order-closed},
%$\mathcal{S} \neq\Proj{\mathcal{S}}$, therefore
$\mathcal{S}$ is not order-closed.
%\begin{obs}
%\label{obs-order-proj-1}
%Let $\mathcal{S}$ be an order-closed set of scenarios, such that all
%its members have only default constraints. Then,
%$\mathcal{S}=\Proj{\mathcal{S}}$.
%\end{obs}
%\noindent

In general, $\ParS{\mathcal{S}}$ captures the ordering relations that are
identical between all the members of $\mathcal{S}$.
If $\mathcal{S}$ is not order-closed, there are certain subtle
similarities between some of the members, similarities that cannot be
explained by $\ParS{\mathcal{S}}$.
For instance, in
Example \ref{ex-one-same-flag}, $\ParS{\mathcal{S}}$ does not
reflect the fact that $\mathtt{flag}$ does not occur immediately after
$\mathtt{start}$ in any member of $\mathcal{S}$: $\mathtt{flag}$ is
preceded by $\mathtt{finish_1}$, $\mathtt{finish_2}$, or both in all
the members of $\mathcal{S}$.

Intuitively, if $\mathcal{S}$ is not order-closed,
%its underlying partial order,
$\ParS{\mathcal{S}}$ is \emph{too permissive}. So
$\ParS{\mathcal{S}}$ would have to be made stricter, somehow, to
reflect \emph{all} the ordering relations that exist between the
members of $\mathcal{S}$. But this cannot be
achieved without partitioning $\mathcal{S}$.
Our goal is to divide $\mathcal{S}$ in such a way that each
partition is order-closed. To achieve this we investigate the
relations between events of $\Sigma_{\mathcal{S}}$ that are not
related by $\ParS{\mathcal{S}}$.
%--------------------
\begin{definition}
\label{def-OBD}
Let $\mathcal{S}$ be a finite set of scenarios over
$\Sigma_{\mathcal{S}}$, such that $\mathcal{S}$ is not order-closed.
Let $e,e'\in\Sigma_{\mathcal{S}}$, such that
$(e,e')\not\in\ParS{\mathcal{S}}$.
The \emph{order-based division} (\emph{OBD} for short)
%The \emph{division}
of $\mathcal{S}$ with respect to $(e,e')$,
denoted by $\OBD{\mathcal{S}}(e,e')$, is
the unordered pair $(\mathcal{S}_1,\mathcal{S}_2)$, where
$\mathcal{S}_1=$\Set{\SeqS_1\in\mathcal{S}}{e\ParS{\SeqS_1}e'} and
$\mathcal{S}_2=$\Set{\SeqS_2\in\mathcal{S}}{e'\ParS{\SeqS_2}e}.
%We say $(\mathcal{S}_1,\mathcal{S}_2)$ is an \emph{order-based
%division} (\emph{OBD} for short) of $\mathcal{S}$.
%\vspace{-0.2em}
\end{definition}
\noindent
Observe that all scenarios in $\mathcal{S}$ contain $e$ and $e'$, but
not all in the same order.
Clearly, an OBD $(\mathcal{S}_1,\mathcal{S}_2)$ is a partitioning,
i.e., $\mathcal{S}_1\neq\emptyset$, $\mathcal{S}_2\neq\emptyset$ and
$\mathcal{S}_1\cup\mathcal{S}_2=\mathcal{S}$.
%---------------
%Notice that
%$\mathcal{S}_1\cup\mathcal{S}_2=\mathcal{S}$,
%$\mathcal{S}_1\neq\emptyset$ and $\mathcal{S}_2\neq\emptyset$:
%all scenarios contain $e$ and $e'$, but not all in the same order.
%\vspace{-0.2em}
\begin{obs}
\label{obs-not-order-closed}
Let $\mathcal{S}$ be a finite set of scenarios, such that
$\mathcal{S}$ is not order-closed and $|\mathcal{S}|>1$.
Then
there is at least one OBD.
\end{obs}
%-------------
\noindent
In Example \ref{ex-one-same-flag}
$\mathtt{finish_1}\not\ParS{\mathcal{S}}\mathtt{finish_2}$ and
$\OBD{\mathcal{S}}(\mathtt{finish_1},\mathtt{finish_2})=(\mathcal{S}_{1},\mathcal{S}_2)$, where $\mathcal{S}_{1}=\{\SeqS_1,\SeqS_2\}$ and
$\mathcal{S}_2=\{\SeqS_3,\SeqS_4\}$.
Notice that $\mathtt{finish_1}$
occurs before $\mathtt{flag}$ in all members of $\mathcal{S}_{1}$,
while in all members of $\mathcal{S}_{2}$
$\mathtt{finish_2}$ occurs before $\mathtt{flag}$.
So

\noindent
%$\ParS{\mathcal{S}}$ with the \emph{additional} ordering relationships:
$\ParS{\mathcal{S}_{1}}=\ParS{\mathcal{S}}\cup\{(\mathtt{finish_1},\mathtt{finish_2}),
(\mathtt{finish_1},\mathtt{flag})\}$ and

\noindent
$\ParS{\mathcal{S}_{2}}=\ParS{\mathcal{S}}\cup\{(\mathtt{finish_2},
\mathtt{finish_1}),(\mathtt{finish_2},\mathtt{flag})\}$.

\Fig{fig-chains-2} shows the maximal chains representing
$\ParS{\mathcal{S}_{1}}$ and $\ParS{\mathcal{S}_2}$.
Observe that both $\mathcal{S}_{1}$ and $\mathcal{S}_{2}$ are
order-closed. Indeed, ${\mathcal{S}}_{1}$ can be realised by
the DTS $\xi_1||\xi_2$, where
$\xi_1=(\mathtt{start}~ \mathtt{finish_1}~ \mathtt{flag},\emptyset)$
and
$\xi_2=(\mathtt{start}~ \mathtt{finish_1}~ \mathtt{finish_2},\emptyset)$.
Similarly, $\mathcal{S}_2$ is realised by %expression
the DTS $\xi_3||\xi_4$, where
$\xi_3=(\mathtt{start}~ \mathtt{finish_2}~ \mathtt{flag},\emptyset)$
and
$\xi_4=(\mathtt{start}~ \mathtt{finish_2}~ \mathtt{finish_1},\emptyset)$.
So
$\mathcal{S}$ is realised by $(\xi_1||\xi_2)\BAR(\xi_3||\xi_4)$.
%expression $E|E'=(E_1||E_2)|(E_3||E_4)$.
%-----------------------

In this example there are two other ways to divide $\mathcal{S}$:

\noindent
$\OBD{\mathcal{S}}(\mathtt{finish_1},\mathtt{flag})=(\{\SeqS_1,\SeqS_2,\SeqS_4\},\{\SeqS_3\})$,
from which we obtain the expression $(\xi_5||\xi_6)\BAR\SeqS_3$,
where
$\xi_5=(\mathtt{start}~ \mathtt{finish_1}~ \mathtt{flag},\emptyset)$
and
$\xi_6=(\mathtt{start}~ \mathtt{finish_2},\emptyset)$.

\noindent
$\OBD{\mathcal{S}}(\mathtt{finish_2},\mathtt{flag})=(\{\SeqS_1,\SeqS_3,\SeqS_4\},\{\SeqS_2\})$,
which results in $(\xi_7||\xi_8)\BAR\SeqS_2$,
where
$\xi_7=(\mathtt{start}~ \mathtt{finish_2}~ \mathtt{flag},\emptyset)$
and
$\xi_8=(\mathtt{start}~ \mathtt{finish_1},\emptyset)$.
%However,
But
such solutions do not align with our goal: our algorithm (see
\Sec{sec-algo})
will produce the first solution.
%result from divisions with singleton partitions, such as these two
%expressions.
%---------------
%In general, ${\mathcal{S}}$ can be partitioned in more than
%one way, so we need some criteria to compare all the possible
%divisions.
%-------------
\begin{definition}
\label{o-ranking}
The \emph{o-ranking} of $(\mathcal{S}_1,\mathcal{S}_2)$,  a division
of $\mathcal{S}$, is
$o_1*\frac{|\mathcal{S}_1|}{|\mathcal{S}|}+o_2*\frac{|\mathcal{S}_2|}{|\mathcal{S}|}$,
where $o_i$ ($i\in\{1,2\}$) is
0 if $\mathcal{S}_i$ is order-closed, and 1 otherwise.
%the number of partitions of $D$ that are not order-closed.
\end{definition}
%---------------
\noindent
This ranking will be used to choose between several possible
divisions: lower o-ranking should give more satisfactory results.
The intuition is that an order-closed partition is
likely to require less additional partitioning.
Moreover, the larger the order-closed partition, the better.
%(notice that singleton partitions are order-closed).
%If two divisions have only one order-closed
%partition, we should favour the one whose order-closed partition is larger.

Observe that in Example \ref{ex-one-same-flag} all the OBDs have equal
o-ranking.
%both partitions of each of the three OBDs are order-closed.
So in order to choose the best
division we must find another criterion that looks more deeply into
the partitions of OBDs.

The following definition is an adaptation of \emph{Kendall tau
distance}\cite{Kendall1938}.
%\vspace{-0.1em}
\begin{definition}
\label{def-K-distance}
Let $\SeqS_1$ and $\SeqS_2$ be two normalised scenarios
of equal length over the same set of events.
The \emph{K-distance} of $\SeqS_1$ and $\SeqS_2$, denoted by
$\K{\SeqS_1}{\SeqS_2}$, is
$|\Set{(e,e')}{e\PR{\SeqS_1}e'\wedge e'\PR{\SeqS_2}e}|$.
\end{definition}
\noindent
The larger the K-distance, the more dissimilar are the sequences of
events.
%of two scenarios.
%\vspace{-0.2em}
\begin{definition}
\label{def-K-set-ranking}
Let $\mathcal{S}$ be a finite set of scenarios.
The \emph{K-ranking} of ${\mathcal{S}}$, denoted by $\RK{\mathcal{S}}$,
is the average of all pairwise K-distances between the members of
$\mathcal{S}$.
\end{definition}
\noindent
Since $\K{\SeqS_1}{\SeqS_2}=\K{\SeqS_2}{\SeqS_1}$,
the average is taken over the number of unordered pairs in
${\mathcal{S}}$, i.e.,
%there are
$\frac{|{\mathcal{S}}|*(|{\mathcal{S}}|-1)}{2}$.
%unordered pairs in ${\mathcal{S}}$.
%the average is taken over the number of unordered pairs. There are
%$\frac{|{\mathcal{S}}|*(|{\mathcal{S}}|-1)}{2}$ such pairs in ${\mathcal{S}}$.
 %----------------
 %\vspace{-0.2em}
\begin{definition}
\label{def-K-ranking}
Let $\mathcal{S}$ be a finite set of scenarios and
$D=(\mathcal{S}_1,\mathcal{S}_2)$ be an OBD of
$\mathcal{S}$. The \emph{K-ranking} of $D$, denoted by
$\RK{D}$, is given by
$\RK{\mathcal{S}_1}*\frac{|\mathcal{S}_1|}{|\mathcal{S}|}+
\RK{\mathcal{S}_2}*\frac{|\mathcal{S}_2|}{|\mathcal{S}|}$.
\end{definition}
\noindent
It turns out that the K-rankings of divisions provide us with
the crucial information that we need in order to choose the best
division when there are multiple divisions that all agree
with each other on their o-ranking.
%-----------------------
%%%%%%%%%%%%% three_in_one %%%%%%%%%
\begin{table}[t]
\begin{center}
\begin{tabular}{| c | c | c | c | c | c|}
\hline
pair & OBD & order-closed &o-rank & K-rank & expression\\
%pairs & divisions & informative? & K-rankings & expressions\\
\hline
$(a,d)$ &
\begin{tabular}{c}
$\{\SeqS_1,\SeqS_3,\SeqS_4,\SeqS_5,\SeqS_6\}$\\
 $\{\SeqS_2\}$
\end{tabular} &
 \begin{tabular}{c}
%$(b,c)$\\
No\\
Yes
\end{tabular} & 0.83 & $1.5$  &
\begin{tabular}{c}
$(\xi_1||\xi_2)\BAR(\xi_3||\xi_4)\BAR\SeqS_2$ \\
 $\xi_1=(abc,\emptyset),\xi_2=(ad,\emptyset)$\\
 $\xi_3=(acb,\emptyset),\xi_4=(acd,\emptyset)$
\end{tabular}
\\
%-----------------------
\hline
 $(b,c)$ & \begin{tabular}{c}
 $\{\SeqS_1,\SeqS_2,\SeqS_5,\SeqS_6\}$\\
 $\{\SeqS_3,\SeqS_4\}$
\end{tabular} &
 \begin{tabular}{c}
% $(a,d),(c,d)$\\
Yes\\
Yes
\end{tabular} &0 &1.44  &
\begin{tabular}{c}
$(\xi_1||\xi_2)\BAR(\xi_3||\xi_4)$ \\
 $\xi_1=(abc,\emptyset),\xi_2=(d,\emptyset)$\\
 $\xi_3=(acb,\emptyset),\xi_4=(acd,\emptyset)$
\end{tabular}
\\
%----------------------
\hline
$(b,d)$ & \begin{tabular}{c}
$\{\SeqS_1,\SeqS_2,\SeqS_3\}$ \\
$\{\SeqS_4,\SeqS_5,\SeqS_6\}$
\end{tabular} &
 \begin{tabular}{c}
% $(c,d)$ \\
No\\
Yes
\end{tabular} &0.5 &1.67 &
\begin{tabular}{c}
$(\xi_1||\xi_2)\BAR\SeqS_3\BAR(\xi_3||\xi_4)$ \\
 $\xi_1=(abc,\emptyset),\xi_2=(dbc,\emptyset)$\\
 $\xi_3=(abd,\emptyset),\xi_4=(ac,\emptyset)$
\end{tabular}
\\
%-------------------
\hline
$(c,d)$ & \begin{tabular}{c}
$\{\SeqS_1,\SeqS_2,\SeqS_6\}$\\
$\{\SeqS_3,\SeqS_4,\SeqS_5\}$
\end{tabular} &
 \begin{tabular}{c}
% $(b,c),(d,b)$\\
%  No
Yes\\
Yes
\end{tabular} &0 &1.34 &
\begin{tabular}{c}
$(\xi_1||\xi_2)\BAR(\xi_3||\xi_4)$ \\
 $\xi_1=(abc,\emptyset),\xi_2=(dc,\emptyset)$\\
 $\xi_3=(ab,\emptyset),\xi_4=(acd,\emptyset)$
\end{tabular}\\
\hline
\end{tabular}
\caption{Divisions, their o-rankings, K-rankings and resulting expressions
 (Example \ref{ex-paper-table})}
\label{example-distance}
\vspace{-4mm}
\end{center}
\end{table}
%----------------------------
%\vspace{-0.2em}
\begin{example}
\label{ex-paper-table}
%Consider the set
Let $\mathcal{S}=\{\SeqS_1,\SeqS_2,\SeqS_3,\SeqS_4,\SeqS_5,\SeqS_6\}$,
%of (sequential) scenarios,
where
$\SeqS_1=(adbc,\emptyset)$,
$\SeqS_2=(dabc,\emptyset)$,
$\SeqS_3=(acdb,\emptyset)$,
$\SeqS_4=(acbd,\emptyset)$,
$\SeqS_5=(abcd,\emptyset)$ and
$\SeqS_6=(abdc,\emptyset)$.
The set $\mathcal{S}$ is not order-closed, so it cannot be realised by
a single DTS.
%The underlying partial order in
The relation $\ParS{\mathcal{S}}$
is represented by the maximal chains $a-b$, $a-c$ and $a-d$.
%, while $d$ is unrelated to all the other events.
There are four pairs of events that are unrelated by
$\ParS{\mathcal{S}}$: $(a,d)$, $(b,c)$,
$(b,d)$ and $(c,d)$. The divisions with respect to these pairs are shown in
the second column of Table \ref{example-distance}, the third
column shows whether each partition of a division is order-closed or
not, while the fourth column shows the o-rankings. Observe that $\OBD{\mathcal{S}}(b,c)$ and
$\OBD{\mathcal{S}}(c,d)$ have equal and lowest o-rankings. So we compute their K-rankings.
%For example, consider the division
% $\OBD{\mathcal{S}}(c,d)=(\{\SeqS_1,\SeqS_2,\SeqS_6\},\{\SeqS_3,\SeqS_4,\SeqS_5\},)$ in the last row of the table.

\noindent
For $\OBD{\mathcal{S}}(c,d)$,
%\noindent
$\RK{\{\SeqS_1,\SeqS_2,\SeqS_6\}}=\frac{\K{\SeqS_1}{\SeqS_2}+\K{\SeqS_1}{\SeqS_6}+\K{\SeqS_2}{\SeqS_6}}{3}\approx
1.33$.
Similarly, $\RK{\{\SeqS_3,\SeqS_4,\SeqS_5\}}\approx
1.33$.
So
%the K-ranking of the OBD is
$\RK{\OBD{\mathcal{S}}(c,d)}=1.33*\frac{3}{6}+1.33*\frac{3}{6}\approx 1.34$.
Similarly, $\RK{\OBD{\mathcal{S}}(b,c)}\approx 1.44$.
%The K-ranking of $\OBD{\mathcal{S}}(b,c)\approx 1.44$.
So we choose $\OBD{\mathcal{S}}(c,d)$ that has a lower K-ranking.
The resulting expression is shown in the last row,
last column of the table:
%. The interpretation is that
$\{\SeqS_1,\SeqS_2,\SeqS_6\}$ is realised by $\xi_1||\xi_2$,
while $\{\SeqS_3,\SeqS_4,\SeqS_5\}$ is realised by $\xi_3||\xi_4$.

The K-rankings of the other two divisions are shown in
the fifth column, while the last column shows the corresponding
expressions.
%All these expressions are equivalent.
Notice that those partitions of $\OBD{\mathcal{S}}(a,d)$ and
$\OBD{\mathcal{S}}(b,d)$
that are not order-closed must be further partitioned.
Observe that $\OBD{\mathcal{S}}(b,c)$ would have also resulted in
an expression with two alternatives. In general, divisions with
lowest K-rankings produce more satisfactory results, unless their
o-rankings differ.
\end{example}
\noindent
In our algorithm (see \Sec{sec-algo}) we
choose a division with the lowest o-ranking.
We break a tie by choosing a division with the lowest K-ranking.
%We show how to break a tie in \Sec{sec-full}.
In order to further break a tie we introduce a new ranking, called l-ranking,
that takes into account whether the distance functions of the
partitions of a division include gaps or not, and how balanced the
partitions are with respect to their sizes.
For a division $D$, its l-ranking is computed by
$100*g+b$, where $g$ is the number of partitions of $D$
that include gaps, and $b$
is the difference between the sizes of the two partitions of $D$.
%Order-closedness is a more indicative factor than
The aim is to give higher priority to divisions whose partitions are
without gaps. Parameter $b$
can be used to discard trivial partitions, in favour of non-trivial
ones. (We assume $| S | < 100$.)


%We break a tie by choosing a division whose partitions are
%order-closed. A division with two order-closed partitions has
%priority over one that has only one order-closed partition, and that
%has priority over one that has no order-closed partitions.
%If order-closedness of partitions cannot break a tie, we randomly
%choose a division, while avoiding partitions with gaps if possible.
%-------------------
%\begin{definition}
%\label{def-informative}
%Let $\mathcal{S}$ be a finite set of scenarios over
%$\Sigma_{\mathcal{S}}$ and
%$(\mathcal{S}_1,\mathcal{S}_2)$ be the OBD
%of $\mathcal{S}$ with respecto to $(e,e')$.
%We say the division is \emph{informative} if
%$\ParS{\mathcal{S}_1}\setminus(\ParS{\mathcal{S}}\cup\{(e,e')\})\neq\emptyset
%~\vee \ParS{\mathcal{S}_2}\setminus(\ParS{\mathcal{S}}\cup\{(e',e)\})\neq\emptyset$.
%\end{definition}
%--------------------
%\noindent
%In Example \ref{ex-one-same-flag} Since
%$(\mathtt{finish_1},\mathtt{flag})\in\ParS{\mathcal{S}_{1}}
%\setminus\ParS{\mathcal{S}}\cup\{(\mathtt{finish_1},\mathtt{finish_2})\}$
%and
%$(\mathtt{finish_2},\mathtt{flag})$\\$\in\ParS{\mathcal{S}_{2}}
%\setminus\ParS{\mathcal{S}}\cup\{(\mathtt{finish_2},\mathtt{finish_1})\}$,
%the division is informative.

%In Example \ref{ex-paper-table}
%All the divisions are informative: the new ordering
%relations are shown in the third column.
